% !TEX program = xelatex
\documentclass[11pt]{article}
\usepackage{fontspec}
\usepackage{xeCJK}
\usepackage[margin=1in]{geometry}
\usepackage{amsmath}
\usepackage{amssymb}
\usepackage{amsthm}
\usepackage{graphicx}
\usepackage{tikz-cd}
\usepackage{multicol}
\usepackage{titling}
\usepackage{setspace}
\usepackage{float}
\usepackage{subcaption}
\usepackage{algorithm}
\usepackage{algpseudocode}
\newtheorem{theorem}{Theorem}
\newtheorem{lemma}{Lemma}

% Bigger & bolder title
\pretitle{\begin{center}\LARGE\bfseries}
\posttitle{\par\end{center}\vskip 0.5em}
\predate{\begin{center}}
\postdate{\end{center}}
\setlength{\droptitle}{-6em}

% Paragraph style (old)
% \setlength{\parindent}{0pt}
% \setlength{\parskip}{0.4\baselineskip}
% Paragraph style (new)
% \setlength{\parindent}{2em}
% \setlength{\parskip}{0pt}
% \setstretch{0.96}


\begin{document}

\title{Notes on ML Models}
\author{}
\date{}
\maketitle

\clearpage
\section{Principal Component Analysis}
PCA algorithm is to find a \textbf{linear and orthogonal} projection of the high dimensional
data $x \in \mathbb{R}^D$, to a low dimensional subspace $z \in \mathbb{R}^L$, such that the low dimensional representation is a "good approximation" to the original data. If we project or encode $mathbf{x}$ to get $\mathbf{z= W^T x}$, and then unproject or decode $\mathbf{z}$ to get $\mathbf{\hat{x}= Wz}$, then we want $\mathbf{x}$ to be close to $\mathbf{x}$ in $\ell_2$ distance. In particular, we can define the following \textbf{reconstruction error}:
\[
\mathcal{L}(\mathbf{W}) \triangleq \frac{1}{N}\sum_{n=1}^{N}
\left\lVert \mathbf{x}_n - \operatorname{decode}(\operatorname{encode}(\mathbf{x}_n;\mathbf{W});\mathbf{W}) \right\rVert_2^2
\]

PCA algorithm states that this objective can be minimized by setting $\mathbf{W} = \mathbf{U}_L$, where $\mathbf{U}_L$ contains the the $L$ eigenvectors with lages eigenvalues of the empirical covariance matrix 
\[
\hat{\boldsymbol{\Sigma}} = \frac{1}{N} \sum_{n=1}^N (\boldsymbol{x}_n - \overline{\boldsymbol{x}})(\boldsymbol{x}_n - \overline{\boldsymbol{x}})^\mathsf{T} = \frac{1}{N} \mathbf{X}_c^\mathsf{T}\mathbf{X}_c
\]
where $\mathbf{X}_c$ is a centered version of the $N \times D$ design matrix.

\subsection{PCA setup}
Suppose we have an (unlabeled) dataset $\mathcal{D} = \{\boldsymbol{x}_n : n = 1:N\}$, where $\boldsymbol{x}_n \in \mathbb{R}^D$. We can represent this as an $N \times D$ data matrix $\mathbf{X}$. We will assume $\overline{\boldsymbol{x}} = \frac{1}{N}\sum_{n=1}^N \boldsymbol{x}_n = \mathbf{0}$, which we can ensure by centering the data.

We would like to approximate each $\boldsymbol{x}_n$ by a low dimensional representation, $\boldsymbol{z}_n \in \mathbb{R}^L$. We assume that each $\boldsymbol{x}_n$ can be explained in terms of a weighted combination of basis functions $\boldsymbol{w}_1, \dots, \boldsymbol{w}_L$, where each $\boldsymbol{w}_k \in \mathbb{R}^D$, and where the weights are given by $\boldsymbol{z}_n \in \mathbb{R}^L$, i.e.
\[
\boldsymbol{x}_n \approx \sum_{k=1}^L z_{nk} \boldsymbol{w}_k
\]
The vector $\boldsymbol{z}_n$ is the low dimensional representation of $\boldsymbol{x}_n$, and is known as the \textbf{latent vector}, the collection of these latent variables are called the \textbf{latent factors}.

The error produced by this approximation:

\[
\mathcal{L}(\mathbf{W}, \mathbf{Z}) = \frac{1}{N} \|\mathbf{X} - \mathbf{Z}\mathbf{W}^\mathsf{T}\|_F^2 = \frac{1}{N} \|\mathbf{X}^\mathsf{T} - \mathbf{W}\mathbf{Z}^\mathsf{T}\|_F^2 = \frac{1}{N} \sum_{n=1}^N \|\boldsymbol{x}_n - \mathbf{W}\boldsymbol{z}_n\|^2
\]
where
\begin{itemize}
    \item $\mathbf{X} \in \mathbb{R}^{N \times D}$: data matrix, rows are centered data points $\mathbf{x}_n^T$
    \item $\mathbf{Z} \in \mathbb{R}^{N \times L}$: latent code matrix, rows are centered data points $\mathbf{z}_n^T$
    \item $\mathbf{W} \in \mathbb{R}^{D \times L}$: basis matrix, columns $\boldsymbol{w}_1, \dots, \boldsymbol{w}_L$ are orthonormal
    \item approximating each $\boldsymbol{x}_n$ by $\hat{\boldsymbol{x}}_n = \mathbf{W}\boldsymbol{z}_n$; or, $\hat{\mathbf{X}
    } = \mathbf{Z} \mathbf{W}^T$.
    \item the minimization of $\mathcal{L}$ is subject to constraints $\mathbf{w}_k^T \mathbf{w_j} = 0 \, \text{for}\, k \neq j \text{;}\, 1 \, \text{for}\, k = j$, i.e., $\mathbf{W}$ is orthogonal
\end{itemize}

\subsection{Algorithm derivation}
\textit{Base case}

Starting by considering the best 1d solution , $\mathbf{w}_1 \in \mathbb{R}^D$, letting the coefficients for each datapoints associated with the first basis vector be denoted by $\tilde{\mathbf{z}}_1 = [z_{11}, \dots , z_{N1}] \in \mathbb{R}^N$. The reconstruction error is 
\[
\begin{aligned}
\mathcal{L}(\boldsymbol{w}_1, \tilde{\boldsymbol{z}}_1) &= \frac{1}{N} \sum_{n=1}^N \|\boldsymbol{x}_n - z_{n1}\boldsymbol{w}_1\|^2 = \frac{1}{N} \sum_{n=1}^N (\boldsymbol{x}_n - z_{n1}\boldsymbol{w}_1)^\mathsf{T}(\boldsymbol{x}_n - z_{n1}\boldsymbol{w}_1) \\
&= \frac{1}{N} \sum_{n=1}^N [\boldsymbol{x}_n^\mathsf{T}\boldsymbol{x}_n - 2z_{n1}\boldsymbol{w}_1^\mathsf{T}\boldsymbol{x}_n + z_{n1}^2 \boldsymbol{w}_1^\mathsf{T}\boldsymbol{w}_1] \\ 
&= \frac{1}{N} \sum_{n=1}^N [\boldsymbol{x}_n^\mathsf{T}\boldsymbol{x}_n - 2z_{n1}\boldsymbol{w}_1^\mathsf{T}\boldsymbol{x}_n + z_{n1}^2]
\quad \text{since} \, \boldsymbol{w}_1^\mathsf{T}\boldsymbol{w}_1 = 1
\end{aligned}
\]

Taking derivatives wrt $z_{n1}$ and equating to zero gives
\[
\frac{\partial}{\partial z_{n1}} \mathcal{L}(\boldsymbol{w}_1, \tilde{\boldsymbol{z}}_1) = \frac{1}{N}[-2\boldsymbol{w}_1^\mathsf{T}\boldsymbol{x}_n + 2z_{n1}] = 0 \Rightarrow z_{n1} = \boldsymbol{w}_1^\mathsf{T}\boldsymbol{x}_n
\]
So the optimal embedding is obtained by orthogonally projecting the data onto $\boldsymbol{w}_1$. Plugging this back in gives the loss for the weights:

\[
\mathcal{L}(\boldsymbol{w}_1) = \mathcal{L}(\boldsymbol{w}_1, \tilde{\boldsymbol{z}}_1^*(\boldsymbol{w}_1)) = \frac{1}{N} \sum_{n=1}^N [\boldsymbol{x}_n^\mathsf{T}\boldsymbol{x}_n - z_{n1}^2] = \text{const} - \frac{1}{N} \sum_{n=1}^N z_{n1}^2
\]
To solve for $\boldsymbol{w}_1$, note that

\[
\mathcal{L}(\boldsymbol{w}_1) = -\frac{1}{N} \sum_{n=1}^N z_{n1}^2 = -\frac{1}{N} \sum_{n=1}^N \boldsymbol{w}_1^\mathsf{T}\boldsymbol{x}_n\boldsymbol{x}_n^\mathsf{T}\boldsymbol{w}_1 = -\boldsymbol{w}_1^\mathsf{T}\hat{\boldsymbol{\Sigma}}\boldsymbol{w}_1
\]
where $\boldsymbol{\Sigma}$ is the empirical covariance matrix (since we assumed the data is centered). We can trivially optimize this by letting $\|\boldsymbol{w}_1\| \to \infty$, so we impose the constraint $\|\boldsymbol{w}_1\| = 1$ and instead optimize

\[
\tilde{\mathcal{L}}(\boldsymbol{w}_1) = \boldsymbol{w}_1^\mathsf{T}\hat{\boldsymbol{\Sigma}}\boldsymbol{w}_1 - \lambda_1(\boldsymbol{w}_1^\mathsf{T}\boldsymbol{w}_1 - 1)
\]
where $\lambda_1$ is a Lagrange multiplier. Taking derivatives and equating to zero we have

\[
\frac{\partial}{\partial \boldsymbol{w}_1} \tilde{\mathcal{L}}(\boldsymbol{w}_1) = 2\hat{\boldsymbol{\Sigma}}\boldsymbol{w}_1 - 2\lambda_1 \boldsymbol{w}_1 = 0
\]

\[
\hat{\boldsymbol{\Sigma}}\boldsymbol{w}_1 = \lambda_1 \boldsymbol{w}_1
\]

Hence the optimal direction onto which we should project the data is an eigenvector of the covariance matrix. Left multiplying by $\boldsymbol{w}_1^\mathsf{T}$ (and using $\boldsymbol{w}_1^\mathsf{T}\boldsymbol{w}_1 = 1$) we find

\[
\boldsymbol{w}_1^\mathsf{T}\hat{\boldsymbol{\Sigma}}\boldsymbol{w}_1 = \lambda_1
\]
Since we want to maximize this quantity (minimize the loss), we pick the eigenvector which corresponds to the largest eigenvalue.

\par\medskip
\noindent \textit{Induction Step}

For another direction $\mathbf{w}_2$, the reconstructon error is
\[
\mathcal{L}(\boldsymbol{w}_1, \tilde{\boldsymbol{z}}_1, \boldsymbol{w}_2, \tilde{\boldsymbol{z}}_2) = \frac{1}{N} \sum_{n=1}^N \|\boldsymbol{x}_n - z_{n1}\boldsymbol{w}_1 - z_{n2}\boldsymbol{w}_2\|^2
\]
Optimizing wrt $\boldsymbol{w}_1$ and $\boldsymbol{z}_1$ gives the same solution as before. $\frac{\partial \mathcal{L}}{\partial \boldsymbol{z}_2} = 0$ yields $z_{n2} = \boldsymbol{w}_2^\mathsf{T}\boldsymbol{x}_n$. Substituting in yields
\[
\mathcal{L}(\boldsymbol{w}_2) = \frac{1}{N} \sum_{n=1}^N [\boldsymbol{x}_n^\mathsf{T}\boldsymbol{x}_n - \boldsymbol{w}_1^\mathsf{T}\boldsymbol{x}_n\boldsymbol{x}_n^\mathsf{T}\boldsymbol{w}_1 - \boldsymbol{w}_2^\mathsf{T}\boldsymbol{x}_n\boldsymbol{x}_n^\mathsf{T}\boldsymbol{w}_2] = \text{const} - \boldsymbol{w}_2^\mathsf{T}\hat{\boldsymbol{\Sigma}}\boldsymbol{w}_2
\]
Dropping the constant term, plugging in the optimal $\boldsymbol{w}_1$ and adding the constraints yields
\[
\tilde{\mathcal{L}}(\boldsymbol{w}_2) = -\boldsymbol{w}_2^\mathsf{T}\hat{\boldsymbol{\Sigma}}\boldsymbol{w}_2 + \lambda_2(\boldsymbol{w}_2^\mathsf{T}\boldsymbol{w}_2 - 1) + \lambda_{12}(\boldsymbol{w}_2^\mathsf{T}\boldsymbol{w}_1 - 0)
\]
with the extra constraint term, it can be shown that the solution is given by the eigenvector with the second largest eigenvalue:
\[
\hat{\boldsymbol{\Sigma}}\boldsymbol{w}_2 = \lambda_2 \boldsymbol{w}_2
\]
The proof continues in this way to show that $\hat{\mathbf{W}} = \mathbf{U}_L$.

\subsection{Computational tricks}
\begin{enumerate}
    \item \textbf{High-dimensional data.} PCA finds the eigenvectors of the $D \times D$ covariance matrix $\mathbf{X}^\mathsf{T}\mathbf{X}$. If $D > N$, it is faster to work with the $N \times N$ matrix $\mathbf{X}\mathbf{X}^\mathsf{T}$. Let $\mathbf{U}$ be an orthonormal matrix containing the eigenvectors of $\mathbf{X}\mathbf{X}^\mathsf{T}$ with corresponding eigenvalues in $\boldsymbol{\Lambda}$. By definition we have $(\mathbf{X}\mathbf{X}^\mathsf{T})\mathbf{U} = \mathbf{U}\boldsymbol{\Lambda}$. Pre-multiplying by $\mathbf{X}^\mathsf{T}$ gives
    \[
    (\mathbf{X}^\mathsf{T}\mathbf{X})(\mathbf{X}^\mathsf{T}\mathbf{U}) = (\mathbf{X}^\mathsf{T}\mathbf{U})\boldsymbol{\Lambda}
    \]
    from which we see that the eigenvectors of $\mathbf{X}^\mathsf{T}\mathbf{X}$ are $\mathbf{V} = \mathbf{X}^\mathsf{T}\mathbf{U}$, with eigenvalues $\boldsymbol{\Lambda}$ as before. However, these eigenvectors are not normalized, since $\|\boldsymbol{v}_j\|^2 = \boldsymbol{u}_j^\mathsf{T}\mathbf{X}\mathbf{X}^\mathsf{T}\boldsymbol{u}_j = \lambda_j \boldsymbol{u}_j^\mathsf{T}\boldsymbol{u}_j = \lambda_j$. The normalized eigenvectors are given by
    \[
    \mathbf{V} = \mathbf{X}^\mathsf{T}\mathbf{U}\boldsymbol{\Lambda}^{-\frac{1}{2}}
    \]
    \item \textbf{Computing PCA using SVD.} Let $\mathbf{U}_\Sigma \boldsymbol{\Lambda}_\Sigma \mathbf{U}_\Sigma^\mathsf{T}$ be the top $L$ eigendecomposition of the covariance matrix $\boldsymbol{\Sigma} \propto \mathbf{X}^\mathsf{T}\mathbf{X}$ (assume $\mathbf{X}$ is centered). Recall the optimal estimate of the projection weights $\mathbf{W}$ is given by the top $L$ eigenvalues, so $\mathbf{W} = \mathbf{U}_\Sigma$. Now let $\mathbf{U}_X \mathbf{S}_X \mathbf{V}_X^\mathsf{T} \approx \mathbf{X}$ be the $L$-truncated SVD approximation to the data matrix $\mathbf{X}$. By SVD definition, the columns of $\mathbf{V}_X$ are the eigenvectors of $\mathbf{X}^\mathsf{T}\mathbf{X}$, so 
    \[
    \mathbf{V}_X = \mathbf{U}_\Sigma = \mathbf{W}.
    \]
    In addition, the eigenvalues of the covariance matrix are related to the singular values of the data matrix via $\lambda_k = s_k^2/N$. Now suppose we are interested in the principal components (or PC scores), rather than the projection matrix. We have
    \[
    \mathbf{Z} = \mathbf{X}\mathbf{W} = \mathbf{U}_X \mathbf{S}_X \mathbf{V}_X^\mathsf{T} \mathbf{V}_X = \mathbf{U}_X \mathbf{S}_X
    \]
    Finally, if we want to approximately reconstruct the data, we have
    \[
    \hat{\mathbf{X}} = \mathbf{Z}\mathbf{W}^\mathsf{T} = \mathbf{U}_X \mathbf{S}_X \mathbf{V}_X^\mathsf{T}
    \]
    This is precisely the same as a \textbf{truncated SVD approximation}.
\end{enumerate}
Thus, PCA can be performed either using an eigendecomposition of $\boldsymbol{\Sigma}$ or an SVD decomposition of $\mathbf{X}$. The latter is often preferable, for computational reasons.

\clearpage
\section{Generalized Linear Models}

A GLM is a conditional version of an exponential family distribution $p(y \mid \theta) = h(y) \exp\left(\theta^\top T(y) - A(\theta)\right)$, in which the natural parameters are a linear function of the input. More precisely, the model has the following form:
\[
p(y_n \mid \mathbf{x}_n, \mathbf{w}, \sigma^2)
= \exp\left[\frac{y_n\eta_n - A(\eta_n)}{\sigma^2} + \log h(y_n, \sigma^2)\right]
\]
where 
\begin{itemize}
    \item $y_n$ is the output/response variable, $\mathbf{x_n}$ is the input/feature vector, $\mathbf{w}$ and $\sigma$ is the learnable model parameter 
    \item $\eta_n \triangleq \mathbf{w}^{\mathsf{T}}\mathbf{x}_n$ is the (input dependent) natural parameter or canonical parameter
    \item $A(\eta_n)$ is the log normalizer
    \item $\mathcal{T}(y_n) = y_n$ is the sufficient statistic, which is the smallest summary of the output data $y$ that retains everything relevant for estimating the parameter $\theta$ -- while it may still lose other information about $y$ itself.
    \item $\sigma^2$ is the dispersion term
    \item $h(y_n, \sigma^2)$  is the scaling constant or base measure
\end{itemize}

We can show that the mean and variance of the response variable are as follows:
\[
\int p(y_n \mid \eta_n, \sigma^2)\, dy_n = 1 \quad \text{for all } \eta_n
\]

To derive the mean, differentiate the equation w.r.t to $\eta_n$
\[ \int \frac{\partial p(y_n \mid \eta_n, \sigma^2)}{\partial \eta_n}\, dy_n = 0 \]
\[\int p(y_n\mid\eta_n,\sigma^2) \cdot \frac{y_n - A'(\eta_n)}{\sigma^2}\, dy_n = 0\]
\[\frac{1}{\sigma^2}\left[\int y_n\, p(y_n\mid\eta_n,\sigma^2)\, dy_n - A'(\eta_n)\right] = 0\]
thus, the \textbf{mean} is:
\[
\mathbb{E}\left[y_n \mid \mathbf{x}_n, \mathbf{w}, \sigma^2\right]
= A'(\eta_n) \triangleq \ell^{-1}(\eta_n)
\]
where the mapping from the linear inputs to the mean of the output using $\mu_n = \ell^{-1}(\eta_n)$, and the function $\ell$ is known as the link function, and $\ell^{-1}$ is known as the mean function. 

To derive the variance, differentiate the above equation w.r.t $\eta_n$ again,
\[
\int \left[-A''(\eta_n)\cdot p + (y_n - A'(\eta_n)) \cdot p \cdot \frac{y_n - A'(\eta_n)}{\sigma^2}\right] dy_n = 0
\]
\[-A''(\eta_n)\int p\, dy_n + \frac{1}{\sigma^2}\int (y_n - A'(\eta_n))^2\, p\, dy_n = 0
\]
thus, the \textbf{variance} is:
\[
\mathbb{V}\left[y_n \mid \mathbf{x}_n, \mathbf{w}, \sigma^2\right]
= A''(\eta_n)\sigma^2
\]

\subsection{Logistic regression}

Binary logistic regression is used when $y_n \in \{0,1\}$. It has the form
\[
p(y_n \mid \mathbf{x}_n, \mathbf{w})
= \mu_n^{y_n}(1-\mu_n)^{1-y_n}
\]
where
\[
\mu_n = \sigma(\eta_n) = \frac{1}{1 + \exp(-\eta_n)}, \quad
\eta_n = \mathbf{w}^{\mathsf{T}}\mathbf{x}_n.
\]
Hence
\[
\log p(y_n \mid \mathbf{x}_n, \mathbf{w})
= y_n\log \mu_n + (1-y_n)\log(1-\mu_n)
\]
Using $\mu_n = \exp(\eta_n)/(1+\exp(\eta_n))$, in GLM form:
\[
\log p(y_n \mid \mathbf{x}_n, \mathbf{w})
= y_n\eta_n - \log(1 + \exp(\eta_n))
\]
We see that $A(\eta_n) = \log(1 + \exp(\eta_n))$ and hence
\[
\mathbb{E}[y_n] = A'(\eta_n) = \frac{1}{1 + \exp(-\eta_n)},\quad \mathbb{V}[y_n] = A''(\eta_n) = \mu_n(1-\mu_n)
\]
Thus the link function is the logit function, $\ell(\mu_n) = \log\left(\frac{\mu_n}{1-\mu_n}\right)$, and the mean function is the sigmoid function, $\ell^{-1}(\eta_n) = \sigma(\eta_n)$.

\subsection{Linear regression}

Linear regression has the form with $y_n \in \mathbb{R}$
\[
p(y_n \mid \mathbf{x}_n, \mathbf{w}, \sigma^2)
= \frac{1}{\sqrt{2\pi\sigma^2}}
\exp\left(-\frac{1}{2\sigma^2}(y_n - \mathbf{w}^{\mathsf{T}}\mathbf{x}_n)^2\right)
\]
Hence
\[
\log p(y_n \mid \mathbf{x}_n, \mathbf{w}, \sigma^2)
= -\frac{1}{2\sigma^2}(y_n - \eta_n)^2 - \frac{1}{2}\log(2\pi\sigma^2)
\]
where $\eta_n = \mathbf{w}^{\mathsf{T}}\mathbf{x}_n$. In GLM form:
\[
\log p(y_n \mid \mathbf{x}_n, \mathbf{w}, \sigma^2)
= \frac{y_n\eta_n - \frac{\eta_n^2}{2}}{\sigma^2}
- \frac{1}{2}\left(\frac{y_n^2}{\sigma^2} + \log(2\pi\sigma^2)\right)
\]
We see that $A(\eta_n) = \eta_n^2/2$ and hence
\[
\mathbb{E}[y_n] = \eta_n = \mathbf{w}^{\mathsf{T}}\mathbf{x}_n,\quad \mathbb{V}[y_n] = \sigma^2 
\]

\subsection{Binomial regression}

If the response variable is the number of successes in $N_n$ trials, $y_n \in \{0,\ldots,N_n\}$, we can use binomial regression, which is defined by
\[
p(y_n \mid \mathbf{x}_n, N_n, \mathbf{w})
= \operatorname{Bin}(y_n \mid \sigma(\mathbf{w}^{\mathsf{T}}\mathbf{x}_n), N_n)
\]
note that binary logistic regression is the special case when $N_n = 1$. 

The log pdf is given by
\[
\log p(y_n \mid \mathbf{x}_n, N_n, \mathbf{w})
= y_n\log \mu_n + (N_n-y_n)\log(1-\mu_n) + \log \binom{N_n}{y_n}
\]
\[
= y_n\log\left(\frac{\mu_n}{1-\mu_n}\right)
+ N_n\log(1-\mu_n) + \log \binom{N_n}{y_n}
\]
where $\mu_n = \sigma(\eta_n)$. To rewrite this in GLM form, let us define
\[
\eta_n \triangleq \log\left(\frac{\mu_n}{1-\mu_n}\right)
= \log\left[\frac{1}{1+e^{-\mathbf{w}^{\mathsf{T}}\mathbf{x}_n}}\frac{1+e^{-\mathbf{w}^{\mathsf{T}}\mathbf{x}_n}}{e^{-\mathbf{w}^{\mathsf{T}}\mathbf{x}_n}}\right]
= \log\left(\frac{1}{e^{-\mathbf{w}^{\mathsf{T}}\mathbf{x}_n}}\right)
= \mathbf{w}^{\mathsf{T}}\mathbf{x}_n
\]
Hence we can write binomial regression in GLM form as follows:
\[
\log p(y_n \mid \mathbf{x}_n, N_n, \mathbf{w})
= y_n\eta_n - A(\eta_n) + h(y_n)
\]
where
\[
h(y_n) = \log \binom{N_n}{y_n}, \quad
A(\eta_n) = -N_n\log(1-\mu_n) = N_n\log(1+e^{\eta_n})
\]
Hence,
\[
\mathbb{E}[y_n] = \frac{dA}{d\eta_n} = \frac{N_n e^{\eta_n}}{1+e^{\eta_n}} = N_n \mu_n, \quad \mathbb{V}[y_n] = \frac{d^2A}{d\eta_n^2} = N_n \mu_n (1-\mu_n)
\]

\subsection{Poisson regression}

If the response variable is an integer count, $y_n \in \{0,1,\ldots\}$, we can use Poisson regression, which is defined by
\[
p(y_n \mid \mathbf{x}_n, \mathbf{w})
= \operatorname{Poi}(y_n \mid \exp(\mathbf{w}^{\mathsf{T}}\mathbf{x}_n))
\]
where
\[
\operatorname{Poi}(y \mid \mu) = e^{-\mu}\frac{\mu^y}{y!}
\]
is the Poisson distribution.

The log pdf is given by
\[
\log p(y_n \mid \mathbf{x}_n, \mathbf{w})
= y_n\log \mu_n - \mu_n - \log(y_n!)
\]
where $\mu_n = \exp(\mathbf{w}^{\mathsf{T}}\mathbf{x}_n)$. Hence in GLM form we have
\[
\log p(y_n \mid \mathbf{x}_n, \mathbf{w})
= y_n\eta_n - A(\eta_n) + h(y_n)
\]
where $\eta_n = \log(\mu_n) = \mathbf{w}^{\mathsf{T}}\mathbf{x}_n$, $A(\eta_n) = \mu_n = e^{\eta_n}$, and $h(y_n) = -\log(y_n!)$. Hence
\[
\mathbb{E}[y_n] = \frac{dA}{d\eta_n} = e^{\eta_n} = \mu_n, \quad
\mathbb{V}[y_n] = \frac{d^2A}{d\eta_n^2} = e^{\eta_n} = \mu_n
\]

\subsection{Canonical and non-canonical link functions}

Link functions connect three quantities:
\begin{itemize}
    \item $\theta_n$ --- the \textbf{natural parameter} of the \textit{exponential family distribution}.
    \item $\mu_n = \mathbb{E}[y_n]$ --- the \textbf{mean} of the output distribution.
    \item $\eta_n = \boldsymbol{w}^\top \boldsymbol{x}_n$ --- the \textbf{linear predictor}, built from the \textit{inputs} and weights.
\end{itemize}

\noindent The two connecting functions:
$$
\theta_n \xrightarrow{\ A'\ } \mu_n \xrightarrow{\ \ell\ } \eta_n
$$
\begin{itemize}
    \item $A'$ (derivative of the log-partition function) always maps $\theta_n \to \mu_n$. This relationship is fixed by the choice of exponential-family distribution (Bernoulli, Gaussian, Poisson, \ldots) and never changes.
    \item $\ell$, the \textbf{link function}, maps $\mu_n \to \eta_n$. This is the function \emph{you choose} when designing the GLM. Its inverse, $\ell^{-1}$, is the \textbf{mean function}: $\mu_n = \ell^{-1}(\eta_n)$ --- the direction actually used to generate predictions from $\eta_n$.
\end{itemize}

\textbf{Canonical link.} Choosing $\ell = (A')^{-1}$ makes the link exactly undo $A'$:
$$
\eta_n = g(\mu_n) = (A')^{-1}(A'(\theta_n)) = \theta_n
$$
so the natural parameter and linear predictor coincide: $\theta_n = \eta_n$.

\textbf{Non-canonical link.} Choosing any $g \ne (A')^{-1}$ gives
$$
\eta_n = g(A'(\theta_n)) \ne \theta_n \quad \text{in general.}
$$
The mean is still recovered correctly via $\mu_n = g^{-1}(\eta_n)$, but the natural parameter $\theta_n$ is no longer equal to the linear predictor $\eta_n$---recovering it requires the separate map $(A')^{-1}(\mu_n)$.

Examples of links:
\begin{itemize}
    \item Logit: canonical link $\theta_n = \ell(\mu_n) = \log\frac{\mu_n}{1-\mu_n}$.
    \item Complementary log-log: non-canonical link $\eta_n = \ell(\mu_n) = \log(-\log(1-\mu_n))$
\end{itemize}

\subsection{Maximum likelihood estimation}

GLMs can be fit the negative log-likelihood as the objective function (ignoring constant term $h$):
\[
\operatorname{NLL}(\mathbf{w})
= -\log p(\mathcal{D} \mid \mathbf{w})
= -\frac{1}{\sigma^2}\sum_{n=1}^{N}\ell_n 
= -\frac{1}{\sigma^2}\sum_{n=1}^{N}[\eta_n y_n - A(\eta_n)]
\]
where $\eta_n = \mathbf{w}^{\mathsf{T}}\mathbf{x}_n$. 

To find the MLE, we must solve $\nabla \text{NNL}(\mathbf{w}) = \mathbf{g(w)} = 0$. For notational simplicity, we will assume $\sigma^2 = 1$. The gradient for a single term is:
\[
\begin{aligned}
\mathbf{g}_n &= 
-\frac{1}{\sigma^2}\frac{\partial \ell_n}{\partial \mathbf{w}}
=-\frac{1}{\sigma^2} \frac{\partial \ell_n}{\partial \eta_n}\frac{\partial \eta_n}{\partial \mathbf{w}} \\
&=-\frac{1}{\sigma^2} \left(y_n - A'(\eta_n)\right)\mathbf{x}_n
=-\frac{1}{\sigma^2} (y_n - \mu_n)\mathbf{x}_n
\end{aligned}
\]
where $\mu_n = f(\mathbf{w^{\mathsf{T}} x})$, and $f$ is the inverse link function that maps from canonical parameters to mean parameters. If we interpret $(y_n - \mu_n)$ as an error signal, then the gradient $-\frac{1}{\sigma^2}\sum_{n=1}^{N}(y_n - \mu_n)\mathbf{x}_n$ weights each input $\mathbf{x}_n$ by its error, and averages the result.

Gradient-based optimization will find a stationary points where $\mathbf{g(w)} = 0$, we need to examine the Hessian matrix to determine whether it is the global optimum. The Hessian is given by 
\[
\begin{aligned}
\mathbf{H} &= \frac{\partial^2}{\partial \boldsymbol{w}\partial \boldsymbol{w}^\top}\text{NLL}(\boldsymbol{w}) \\ 
&= -\sum_{n=1}^N \frac{\partial \boldsymbol{g}_n}{\partial \boldsymbol{w}^\top}
= -\sum_{n=1}^N \frac{\partial \boldsymbol{g}_n}{\partial \mu_n}\frac{\partial \mu_n}{\partial \boldsymbol{w}^\top}
= -\sum_{n=1}^N \boldsymbol{x}_n f'(\boldsymbol{w}^\top \boldsymbol{x}_n)\boldsymbol{x}_n^\top
\end{aligned} 
\]
In general, the Hessian is positive definite, since $f'(\boldsymbol{w}^\top \boldsymbol{x}_n) > 0$, hence the negative log-likelihood is convex, so the MLE for a GLM is unique given $f(\boldsymbol{w}^\top \boldsymbol{x}_n) > 0$ for all $n$.

Based on these results, gradient-based optimizers can be used to estimate the model weights.

\subsubsection{Stochastic gradient descent}

The goal is to solve
\[
\hat{\boldsymbol{w}} \triangleq \operatorname*{argmin}_{\boldsymbol{w}} \mathcal{L}(\boldsymbol{w}) = \operatorname*{argmin}_{\boldsymbol{w}} \text{NLL}({\mathbf{w}})
\]
One of the methods is stochastic gradient descent: we use a minibatch from the whole dataset of size $N$, then we get the following update equation:
\[
\mathbf{w}_{t+1}
= \mathbf{w}_t - \alpha_t \nabla_{\mathbf{w}}\operatorname{NLL}(\mathbf{w}_t)
= \mathbf{w}_t - \alpha_t(\mu_n - y_n)\mathbf{x}_n
\]
where we replaced the average over all $N$ samples with a stochastically chosen minibatch size at each step $t$.

Since we know the objective is convex, then it can be shown that this will converge to a global optimum given the learning rate $\alpha_t$ is decayed appropriately at each step $t$.

\end{document}
